\documentclass[10pt,a4paper]{amsart}
\usepackage[margin=19mm]{geometry}
\usepackage{amsmath,amssymb,mathtools}
\usepackage[scr=rsfs,cal=euler]{mathalfa}
\usepackage{xfrac}
\usepackage[unicode,hidelinks,bookmarks=false]{hyperref}
\newcommand{\ie}{i.e.\ }
\newcommand{\cf}{cf.\ }
\newcommand{\resp}{respectively\ }
\newcommand{\Subsection}{\subsection}
\providecommand{\nobreakdash}{\nobreak\hbox{-}}
\setlength{\emergencystretch}{2em}
\multlinegap0pt
\usepackage{bm}
\usepackage{bbm}
\usepackage{dsfont}
\usepackage{xspace}
\usepackage{cancel}
\usepackage{mathtools}
\usepackage{amsthm}
\makeatletter
\@ifundefined{thm}{\newtheorem{thm}{Thm}}{}
\@ifundefined{claim}{\newtheorem{claim}{Claim}}{}
\@ifundefined{defn}{\newtheorem{defn}[thm]{Definition}}{}
\@ifundefined{fact}{\newtheorem{fact}[thm]{Fact}}{}
\makeatother
\makeatletter
\expandafter\ifx\csname cproof\endcsname\relax
\newenvironment{cproof}{\paragraph{\emph{Proof.\,}}}{\hfill{$\boxminus$}\par\vspace{2mm}}
\fi
\renewcommand{\mid}{\mathrel{|}\allowbreak}
\makeatother
\title[On Borel subsets of generalized Baire spaces]{On Borel subsets of generalized Baire spaces}
\author{Tapani Hyttinen, Miguel Moreno, Jouko V\"a\"an\"anen}
\date{Source: arXiv:2507.15427v2 (CC BY 4.0)}
\begin{document}
\maketitle

\noindent\emph{Source and licence.} On Borel subsets of generalized Baire spaces. Version arXiv:2507.15427v2 (CC BY 4.0), distributed under the Creative Commons Attribution 4.0 International licence, as recorded in the arXiv licence metadata for this exact version and shown on the abstract page https://arxiv.org/abs/2507.15427v2. Full rights notice: licenses/arxiv-2507.15427-CC-BY-4.0.txt.\par\smallskip

\noindent\textit{Editorial scope.} Counted unit: the Fact whose source label is \texttt{None} in the article (source0.tex lines 1139--1142, source label \texttt{None}), delivered together with the article's own complete original proof of it (source0.tex lines 1144--1160, one \texttt{proof} environment of 17 lines). No mathematics is written, repaired or re-derived: statement and proof are the article's own text line for line. Required context retained uncounted: good\_labelled\_tree\_empt (lines 946--962); good\_labelled\_tree (lines 966--978). Every definition, display and label of those lines is retained unchanged. Omitted from this package and not redistributed: 1--945, 963--965, 979--1138, 1143--1143, 1161--1433. Nothing inside a retained excerpt is omitted or altered: every retained body is delivered exactly as the article's lines are written, so no omission receipt of any kind accompanies this package. Citations: the retained excerpts carry no citation command at all, so no bibliography entry of the article is transcribed and no reference is redistributed. Reference adaptation: none was needed and none was made; every reference inside the delivered unit resolves inside it, so no unresolved reference or citation is produced. Printed slips kept literal: no printed word, symbol, hypothesis or number of the counted statement or of its proof was altered. Layout and class: the article's own class is \texttt{amsart} with packages including amssymb, cmap, verbatim, stmaryrd, hyperref, mathrsfs; the wrapper is the lane's pinned \texttt{amsart} editorial wrapper at 10pt on A4, which restates the theorem-like environments the retained text needs and adds the article's own macro definitions that the retained text needs (\texttt{\textbackslash mid}), with extra wrapper packages (bm, bbm, dsfont, xspace, cancel, mathtools). The delivered numbering is the unit's own. Assets: no retained span contains a figure environment, a graphics-inclusion command, a PDF-page inclusion, a table or a TikZ picture; no image file is copied into this package, the package declares no asset dependency and no image file is redistributed with it. Licence: version arXiv:2507.15427v2 of this article is distributed under the Creative Commons Attribution 4.0 International licence, as recorded in the arXiv licence metadata for this exact version and shown on the abstract page https://arxiv.org/abs/2507.15427v2. A full rights notice is delivered at licenses/arxiv-2507.15427-CC-BY-4.0.txt. Characters: every retained excerpt is pure ASCII, so the delivered engine, XeLaTeX with its default Latin Modern text font, reports no missing character.
\begin{defn}\label{good_labelled_tree_empt}
    We say that $(T,a_t,\mathcal{A}_t)_{t\in T}$ is a good labelled tree of $M$ (we drop $M$ when it is clear from the context) if the following holds:
    \begin{enumerate}
        \item $T$ is a subtree of $(\kappa^+)^{<\omega}$.
        \item If $t^\frown i\in T$, and $j<i$, then $t^\frown j\in T$.
        \item If $ht_T(t)=0$, then $\mathcal{A}_{t}\prec M$ is primary over the empty set and $a_t\in\mathcal{A}_{t}$.
        %\item If $t=\{(0,0)\}$, then $a_t=a_\emptyset=\emptyset$. 
        \item If $t\in T$ and $ht_T(t)\ge 1$, then $\mathcal{A}_t =\mathcal{A}_{t^-}[a_t]\preceq M$ and $a_t\not \subseteq\mathcal{A}_{t^-}$.
       % \item If $t\in T$ and $ht_T(t)\ge 1$, then $tp(a_t / \mathcal{A}_{t^-})$ is non-algebraic.
        \item If $t\in T$ and $ht_T(t)\ge 2$, then $tp(a_t / \mathcal{A}_{t^-})\dashv \mathcal{A}_{t^{--}}$.
        \item If $t^\frown i\in T$, then $$a_{t^\frown i}\downarrow_{\mathcal{A}_t} \bigcup_{j<i}a_{t^\frown j}.$$
    \end{enumerate}
    We say that $(T,a_t,\mathcal{A}_t)_{t\in T}$ is a good labelled tree over a countable $\mathcal{A}\prec M$ of $M$, if $(T,a_t,\mathcal{A}_t)_{t\in T}$ satisfies (1), (2), (4), (5), and (6), and the following:
\begin{itemize}
    \item [($3'$)] If $ht_T(t)=0$, then $\mathcal{A}_{t}=\mathcal{A}$ and $a_t\in \mathcal{A}$.
\end{itemize}
\end{defn}

\begin{defn}\label{good_labelled_tree}
    We say that $(T,a_t,\mathcal{A}_t)_{t\in T}$ is a good labelled tree over $(\mathcal{B},\mathcal{A})$, $\mathcal{B}\prec\mathcal{A}\prec M$ countable, of $M$ if the following holds:
    \begin{enumerate}
        \item $T$ is a subtree of $(\kappa^+)^{<\omega}$.
        \item If $t^\frown i\in T$, $ht_T(t)\ge 1$, and $j<i$, then $t^\frown j\in T$.
        \item If $ht_T(t)=0$, then $\mathcal{A}_{t}=\mathcal{B}$, $t$ has a unique immediate successor, $t^+$, $\mathcal{A}_{t^+}=\mathcal{A}$, $a_t\in \mathcal{B}$ and $a_{t^+}\in \mathcal{A}$.
        %\item If $t=\{(0,0)\}$, then $a_t=a_\emptyset=\emptyset$. 
        \item If $t\in T$ and $ht_T(t)\ge 1$, then $\mathcal{A}_t =\mathcal{A}_{t^-}[a_t]\preceq M$ and $a_t\not \subseteq\mathcal{A}_{t^-}$.
       % \item If $t\in T$ and $ht_T(t)\ge 1$, then $tp(a_t / \mathcal{A}_{t^-})$ is non-algebraic.
        \item If $t\in T$ and $ht_T(t)\ge 2$, then $tp(a_t / \mathcal{A}_{t^-})\dashv \mathcal{A}_{t^{--}}$.% and $tp(a_t / \mathcal{A}_{t^-})$ is non-algebraic.
        \item If $t^\frown i\in T$, then $$a_{t^\frown i}\downarrow_{\mathcal{A}_t} \bigcup_{j<i}a_{t^\frown j}.$$
    \end{enumerate}
\end{defn}

\begin{fact}
    For all $t\in T$ and $t^\frown i\in T$, %Definition \ref{good_labelled_tree} item 6, i.e. $a_{t^\frown i}\downarrow_{\mathcal{A}_t} \bigcup_{j<i}a_{t^\frown j}$, implies 
    $\mathcal{A}_{t^\frown i}\downarrow_{\mathcal{A}_t} \bigcup\{\mathcal{A}_{t^\frown j}\mid j\neq i \text{ and } t^\frown j\in T\}$.
\end{fact}

\begin{proof}
    Let $t\in T$. Due to the finite character, it is enough to show that for all finite $B\subseteq \{j<\kappa^+\mid t^\frown j\in T\}$ and $i\in B$ $$\mathcal{A}_{t^\frown i}\downarrow_{\mathcal{A}_t} \bigcup_{j\in B\backslash \{i\}}\mathcal{A}_{t^\frown j}.$$ Let us proceed by induction on the size of $B$.

    Let $B=\{i,j\}$. By Definition \ref{good_labelled_tree_empt} item 7, $a_{t^\frown i}\downarrow_{\mathcal{A}_t} a_{t^\frown j}$ and by domination $\mathcal{A}_{t^\frown i}\downarrow_{\mathcal{A}_t} \mathcal{A}_{t^\frown j}$.

    Let $n$ be such that for all $B\subseteq \{j<\kappa^+\mid t^\frown j\in T\}$ of size $n$ and $i\in B$ $$\mathcal{A}_{t^\frown i}\downarrow_{\mathcal{A}_t} \bigcup_{j\in B\backslash \{i\}}\mathcal{A}_{t^\frown j}.$$ Let $B\subseteq \{j<\kappa^+\mid t^\frown j\in T\}$ of size $n+1$. Let $\langle b_k\rangle_{k<n+1}$ be an enumeration of $\{a_{t^\frown j}|j\in B\}$ such that if $m<k$ and $b_m=a_{t^\frown h}$ and $b_k=a_{t^\frown s}$, then $h<s$. Let us denote by $\mathcal{A}_k$ the model $\mathcal{A}_{t^\frown j}$, where $b_k=a_{t^\frown j}$.

    \begin{claim}
        For all $0<k<n+1$ $$\bigcup_{j<n+1-k}b_j\downarrow_{\mathcal{A}_t}\bigcup_{n+1-k\leq j<n+1}\mathcal{A}_j.$$
    \end{claim}
    \begin{proof}
        For the case $k=1$, from Definition \ref{good_labelled_tree} item 7, $$\bigcup_{j<n}b_j\downarrow_{\mathcal{A}_t}b_n,$$ since $b_n\triangleright_{\mathcal{A}_{t}}\mathcal{A}_n$, $$\bigcup_{j<n}b_j\downarrow_{\mathcal{A}_t}\mathcal{A}_n.$$

        Let $k<n$ be such that $$\bigcup_{j<n+1-k}b_j\downarrow_{\mathcal{A}_t}\bigcup_{n+1-k\leq j<n+1}\mathcal{A}_j.$$ Therefore $$b_{n+1-(k+1)}\downarrow_{\mathcal{A}_t\cup \bigcup_{j<n+1-(k+1)}b_j}\bigcup_{n+1-k\leq j<n+1}\mathcal{A}_j.$$ From Definition \ref{good_labelled_tree} item 7, $$b_{n+1-(k+1)}\downarrow_{\mathcal{A}_t} \bigcup_{j<n+1-(k+1)}b_j$$ by transitivity $$b_{n+1-(k+1)}\downarrow_{\mathcal{A}_t}\bigcup_{n+1-k\leq j<n+1}\mathcal{A}_j\cup \bigcup_{j<n+1-(k+1)}b_j.$$ By domination $$\mathcal{A}_{n+1-(k+1)}\downarrow_{\mathcal{A}_t}\bigcup_{n+1-k\leq j<n+1}\mathcal{A}_j\cup \bigcup_{j<n+1-(k+1)}b_j,$$ so $$\mathcal{A}_{n+1-(k+1)}\downarrow_{\mathcal{A}_t\cup \bigcup_{n+1-k\leq j<n+1}\mathcal{A}_j} \bigcup_{j<n+1-(k+1)}b_j.$$ Since $$\bigcup_{j<n+1-k}b_j\downarrow_{\mathcal{A}_t}\bigcup_{n+1-k\leq j<n+1}\mathcal{A}_j,$$ by transitivity $$\bigcup_{n+1-(k+1)\leq j<n+1}\mathcal{A}_j\downarrow_{\mathcal{A}_t} \bigcup_{j<n+1-(k+1)}b_j.$$ The claim follows from symmetry.
    \end{proof}
    From the previous claim we know $$b_0\downarrow_{\mathcal{A}_t}\bigcup_{1\leq j<n+1}\mathcal{A}_j$$ and by domination $$\mathcal{A}_0\downarrow_{\mathcal{A}_t}\bigcup_{1\leq j<n+1}\mathcal{A}_j.$$ Finally, let $i\in B$ and $b_k=a_{t^\frown i}$. If $k=0$, then we are done. Let us take care of the case $0<k$. Let $b_0=a_{t^\frown l}$, notice that $B\backslash \{l\}$ has size $n$. By the induction hypothesis, $$\mathcal{A}_{t^\frown i}\downarrow_{\mathcal{A}_t} \bigcup_{j\in B\backslash \{i,l\}}\mathcal{A}_{t^\frown j}.$$ So $$\mathcal{A}_{k}\downarrow_{\mathcal{A}_t} \bigcup_{j<n+1, j\notin \{0,k\}}\mathcal{A}_{j}.$$ Since $$\mathcal{A}_{k}\downarrow_{\mathcal{A}_t\cup \bigcup_{j<n+1, j\notin \{0,k\}}\mathcal{A}_{j}}\mathcal{A}_0,$$ by transitivity $$\mathcal{A}_{k}\downarrow_{\mathcal{A}_t} \bigcup_{j<n+1, j\neq k}\mathcal{A}_{j},$$ $$\mathcal{A}_{t^\frown i}\downarrow_{\mathcal{A}_t} \bigcup_{j\in B\backslash \{i\}}\mathcal{A}_{t^\frown j}.$$
\end{proof}

\end{document}
